\section{From transition measurements to opening equity}
\label{sec:insurance-transition-balances}

The fair-value approach produces a margin or a loss component, but those two
outcomes enter the liability differently. Consider two distinct groups of
issued contracts for which the grouping and eligibility for this transition
method have already been established. Use the historical provisions just
described and assume that both groups contain only remaining coverage. All
amounts below are authored currency units at the same transition date. There
are no incurred claims, reinsurance, acquisition balances, tax, accumulated
OCI or other transition adjustments in this example. These simplifications
isolate the effect of the two measurements; they do not supply actual insurer
valuations or evidence of eligibility.

For group A, suppose the measured fulfilment cash flows are 90 and the
separately measured fair value is 110. The FCF amount already includes the
discounting and nonfinancial risk adjustment required by the measurement
rules. Paragraph C20 makes the difference, $110-90=20$, the CSM. This is
unearned profit associated with future services. Keeping it within the
liability gives $90+20=110$; recording it as immediate profit would contradict
that role \citep[paragraphs 32--38, 40(a) and C20]{aasb17}.

Now reverse the ordering of the inputs. For group B, let FCF be 120 and fair
value 100. The difference identifies a loss component of 20, while CSM is
zero. The liability is 120. Paragraphs 49--50 describe the loss component as
part of the liability for remaining coverage, with its own subsequent
allocation and reversal rules. Adding another 20 to the FCF balance would
count that component twice. Replacing the liability by fair value, 100,
would instead understate the required FCF balance
\citep[paragraphs 40, 47, 49--50 and C20]{aasb17}.

Write $F$ for measured future-service FCF, $V$ for the group's fair value,
$C$ for CSM and $K$ for the nonnegative loss component. Positive amounts
denote liabilities. The two sign cases give the following local derivation:
\[
 C=\max(V-F,0),\qquad K=\max(F-V,0),\qquad
 L=F+C=\max(F,V).
\]
Here $L$ is the recognised remaining-coverage liability. The loss component
$K$ is tracked within $L$ and is not added to it. Thus $C-K=V-F$, and at
$V=F$ both components are zero. These relations apply to the stated
transition case; later changes in FCF, service provision and finance effects
follow the subsequent-measurement rules.

The numerical comparison makes the distinction visible:
\begin{center}
\begin{tabular}{lrr}
\toprule
Currency units at transition & Group A & Group B \\
\midrule
Fulfilment cash flows & 90 & 120 \\
Fair value of the group & 110 & 100 \\
CSM & 20 & 0 \\
Loss component within the liability & 0 & 20 \\
Recognised liability & 110 & 120 \\
\bottomrule
\end{tabular}
\end{center}
The recognised liability equals fair value for A and exceeds it for B.
Calling the method fair-value transition therefore does not mean that every
resulting liability is carried at fair value.

\clearpage
To see the opening-equity effect, suppose the old carrying liabilities are
105 for A and 115 for B. Assets of 500 remain unchanged, and these are the
only liabilities. Before transition, equity is $500-105-115=280$. Afterwards,
it is $500-110-120=270$. Each group's liability increases by 5, so total
opening equity falls by 10. Paragraph C5 replaces the retrospective
measurement requirement in C4(a); the removal of incompatible old balances
and the net adjustment to equity in C4(b)--(c) still apply. The ordinary
initial-recognition loss rule does not turn this transition adjustment into
current-period operating expense \citep[paragraphs C4--C5]{aasb17}.

Illustrative component entries show the same reconciliation. For A, debit
the old liability 105 and opening equity 5, then credit the FCF liability 90
and the CSM liability 20. For B, debit the old liability 115 and opening
equity 5, then credit the FCF liability 120. Record B's loss component of 20
within that 120 for its subsequent allocation; it produces no additional
credit. Neither entry changes cash. An actual transition would also need the
acquisition, reinsurance, tax and OCI effects excluded here.

The order of calculation matters as much as the signs. If an implementation
first combined the two already distinct groups, it would find total FCF of
210 and total fair value of 210. Applying the sign rule to those totals would
produce neither CSM nor a loss component and would report a liability of 210.
The correct group-by-group calculation gives 230. Pooling would therefore
understate liability and overstate equity by 20. The programme must retain
the admitted group boundary while calculating CSM and loss components, then
aggregate the resulting balances. Whether different grouping is itself
permitted requires the separate grouping rules; an aggregation shortcut
cannot decide that accounting question.

The example supplies the arithmetic once the group's two measurements and
old balances are known. A real transition disclosure may reveal only some
of those quantities. The next case shows both what can be read from an
insurer's report and what still requires contract and valuation records.

\clearpage
